% ============================================================
%  template/summary-template.tex — 한 장 정리 (2쪽) 골격
%  ★ 시험 문장 번역기는 반드시 한 tabular{m{6.5cm} c m{4.4cm}} 안에 모든 줄을 넣는다 (규칙 ⑱).
%    줄마다 tabular를 따로 만들면 ⇒ 열이 어긋난다.
%  \reflex, \bnum 은 preamble.tex에 이미 정의돼 있다 (여기서 \newcommand 하지 말 것).
% ============================================================
\newpage
\renewcommand{\bandtext}{한 장 정리 · 시험 직전엔 이 두 페이지만 넘겨라}
\renewcommand{\bandcolor}{prime}
\renewcommand{\pageans}{}

{\sffamily\cjkblack\fontsize{16}{19}\selectfont\textcolor{prime}{한 장 정리}\quad \fontsize{12}{15}\selectfont 단원 이름 (N강--M강)}\quad
{\sffamily\footnotesize\textcolor{gray60}{\cjkmed 빈칸을 먼저 채우고, 막히면 그 강으로}}
\par\vspace{2mm}{\color{prime}\hrule height 1.2pt}\vspace{7mm}

{\sffamily\cjkbold\fontsize{11.5}{14}\selectfont\textcolor{prime}{조건 반사} \; \metabar{문제의 신호를 보면 몸이 먼저 움직여야 한다. 빈칸을 채워 보자. (답은 이 페이지 맨 아래)}}
\par\vspace{4mm}

\reflex{신호 문장을 보면}{\bnum{①}\underline{\hspace{16mm}}\,을 먼저 한다}{N강}
\reflex{신호 문장을 보면}{\bnum{②}\underline{\hspace{16mm}}\,부터}{N강}
\reflex{신호 문장을 보면}{\bnum{③}\underline{\hspace{16mm}}\,로 바꾼다}{M강}

\vfill
{\sffamily\cjkbold\fontsize{11.5}{14}\selectfont\textcolor{prime}{시험 문장 번역기} \; \metabar{출제자의 여러 문장이 사실은 같은 조건이다}}
\par\vspace{3mm}

\begin{tcolorbox}[colback=prime06, colframe=prime30, boxrule=0.6pt, arc=1.5mm,
  left=4mm, right=4mm, top=3.5mm, bottom=3.5mm]
\noindent
\begin{tabular}{@{}m{6.5cm} c m{4.4cm}@{}}
{\small $\left.\begin{array}{@{}l@{}}\mbox{``문장 A''}\\[3pt] \mbox{``문장 A$'$''}\end{array}\right\}$} & {\large$\Rightarrow$} & {\small \cjkbold 같은 조건} \\[5mm]
{\small ``문장 B''} & {\large$\Rightarrow$} & {\small 조건 B} \\[4.5mm]
{\small ``문장 C''} & {\large$\Rightarrow$} & {\small 조건 C} \\
\end{tabular}
\end{tcolorbox}

\vfill
\noindent
{\sffamily\fontsize{8.5}{11}\selectfont\cjkmed\textcolor{gray60}{조건 반사 정답 \textcolor{gray30}{|}\; ① 답 \;·\; ② 답 \;·\; ③ 답}}

\newpage
\renewcommand{\pageans}{}

{\sffamily\cjkbold\fontsize{11.5}{14}\selectfont\textcolor{prime}{판정 흐름도} \; \metabar{이 단원의 문제는 전부 이 지도 위에 있다}}
\par\vspace{4mm}

\begin{center}
\begin{tikzpicture}[
  box/.style={fill=prime12, rounded corners=1.2mm, inner xsep=3mm, inner ysep=2mm, font=\sffamily\small},
  navybox/.style={fill=prime, text=white, rounded corners=1.2mm, inner xsep=3mm, inner ysep=2mm, font=\sffamily\small},
  lightbox/.style={draw=prime30, fill=white, rounded corners=1.2mm, inner xsep=2.5mm, inner ysep=1.8mm, font=\small},
  arr/.style={->, prime70, thick}]
  \node[navybox] (start) at (0,0) {\cjkbold 시작};
  \node[box] (s1) at (-3.5,-1.5) {\cjkmed 갈래 1};
  \node[box] (s2) at (3.5,-1.5) {\cjkmed 갈래 2};
  \node[lightbox, anchor=west] (f1) at (-6.2,-3.0) {갈래 1 → 처방};
  \node[lightbox, anchor=west] (f2) at (-6.2,-4.2) {갈래 2 → 처방};
  \draw[arr] (start) -- (s1); \draw[arr] (start) -- (s2);
  \draw[arr] (s1.south) -- ++(0,-0.6) -| (-6.6,-3.0) -- (f1.west);
  \draw[arr] (s2.south) -- ++(0,-0.6) -| (-6.6,-4.2) -- (f2.west);
\end{tikzpicture}
\end{center}

\vspace{5mm}
{\sffamily\cjkbold\fontsize{11.5}{14}\selectfont\textcolor{prime}{실수 클리닉} \; \metabar{O/X로 점검. 하나라도 헷갈리면 해당 위치로}}
\par\vspace{3mm}

\begin{tcolorbox}[colback=prime06, colframe=prime30, boxrule=0.6pt, arc=1.5mm,
  left=4mm, right=4mm, top=3.5mm, bottom=3.5mm]
\begin{itemize}[leftmargin=5mm, itemsep=4pt, topsep=1pt]
\item[] {\cjkbold(\,\;\,)} 명제 1. \hfill {\footnotesize\textcolor{gray60}{X. 이유. 오진 주의 N}}
\item[] {\cjkbold(\,\;\,)} 명제 2. \hfill {\footnotesize\textcolor{gray60}{O. N강 개념 ①}}
\end{itemize}
\end{tcolorbox}
